\section{Алгебры с одной операцией. Полугруппы. Моноиды. Подстановки, композиция подстановок. Две теоремы. Примеры.}

\subsection*{Полугруппы}

\paragraph{Определение}
\textbf{Полугруппа} — это алгебра $\mathcal{S} = (M; \ast)$ с одной ассоциативной бинарной операцией:
$$\forall a, b, c \in M: a \ast (b \ast c) = (a \ast b) \ast c.$$
Если в полугруппе существует система образующих из одного элемента, она называется \textbf{циклической} (пример: $(\mathbb{N}; +)$ с образующей $\{1\}$).

\paragraph{Определяющие соотношения}
Если элементы полугруппы представляются как слова в алфавите образующих $A$, то равенства слов $\alpha = \beta$, определяющие один и тот же элемент, называются \textbf{определяющими соотношениями}.
\begin{itemize}
    \item \textbf{Теорема (Маркова — Поста):} Существует полугруппа, в которой проблема распознавания равенства слов \textbf{алгоритмически неразрешима}.
\end{itemize}

\subsection*{Моноиды и подстановки}

\paragraph{Определение моноида}
\textbf{Моноид} — это полугруппа с нейтральным элементом (единицей) $e$:
$$\exists e \in M \ (\forall a \in M: a \ast e = e \ast a = a).$$

\paragraph{Подстановки (в смысле замены переменных)}
По Новикову, \textbf{подстановкой} называется множество пар $\sigma = \{t_i // v_i\}$, где $v_i$ — переменные, а $t_i$ — термы.
\begin{itemize}
    \item \textbf{Композиция подстановок} $\sigma_1 \circ \sigma_2$ — это новая подстановка, полученная последовательным применением замен. Множество подстановок образует моноид, где единица — тождественная подстановка $\{v_i // v_i\}$.
\end{itemize}

\subsection*{Две теоремы о моноидах}

\paragraph{Теорема 1 (Единственность единицы)}
Единица моноида единственна.

\textit{Доказательство:}
Пусть существуют две единицы $e_1$ и $e_2$. Тогда:
\begin{enumerate}
    \item $e_1 \ast e_2 = e_1$ (так как $e_2$ — единица).
    \item $e_1 \ast e_2 = e_2$ (так как $e_1$ — единица).
\end{enumerate}
Из (1) и (2) следует, что $e_1 = e_2$. $\square$

\paragraph{Теорема 2 (О представлении)}
Всякий моноид $\mathcal{M}$ над множеством $M$ изоморфен некоторому моноиду преобразований над $M$.

\textit{Доказательство:}
Пусть $\mathcal{M} = (M; \ast)$ — моноид. Построим моноид преобразований $\mathcal{M}' = (F; \circ)$, где $F := \{f_m: M \to M \mid f_m(x) := x \ast m\}$.
\begin{enumerate}
    \item $\mathcal{M}'$ — моноид, так как суперпозиция функций ассоциативна, $f_e$ — тождественная функция ($f_e(x) = x \ast e = x$), и $F$ замкнуто относительно $\circ$:
    $$f_{a \ast b}(x) = x \ast (a \ast b) = (x \ast a) \ast b = f_a(x) \ast b = f_b(f_a(x)) = (f_a \circ f_b)(x).$$
    \item Построим отображение $h: M \to F$, где $h(m) := f_m$.
    \item $h$ — гомоморфизм: $h(a \ast b) = f_{a \ast b} = f_a \circ f_b = h(a) \circ h(b)$.
    \item $h$ — биекция:
    \begin{itemize}
        \item Сюръекция по построению множества $F$.
        \item Инъекция: так как $f_a(e) = e \ast a = a$ и $f_b(e) = e \ast b = b$, то из $a \neq b$ следует $f_a \neq f_b$.
    \end{itemize}
\end{enumerate}
Следовательно, $h$ — изоморфизм. $\square$

\subsection*{Примеры}
\begin{enumerate}
    \item $(\mathbb{N}_0; +)$ — моноид натуральных чисел с нулем.
    \item $A^*$ — моноид слов в алфавите $A$ (включая пустое слово $e$).
    \item $(\mathbb{Z}; \cdot)$ — моноид целых чисел по умножению (единица — $1$).
    \item Множество всех функций $f: M \to M$ относительно суперпозиции.
\end{enumerate}